\section{从电流模型改变到材料信息、残差方向和实际收益}
\subsection{为什么要从一个可靠的接收基底出发}
接收谱基底（receiver spectral backbone）按电流直接产生的接收响应保留方向。它的好处是定义清楚、容易复现，而且已有比较显示它在较大的电流预算下相当可靠。现在的问题不是从空集重新寻找所有方向，而是保留这一基础，只检查少量替换是否值得。

一次替换有两个端点：原空间 $[V,A]$ 和新空间 $[V,C]$。$V$ 是共同保留的电流核心（common retained core）；$A$ 是移出的方向，$C$ 是移入的方向。两端使用相同材料状态、照明、正则化和测量残差。若共同核心和相应Schur块可逆，前面的反馈分解给出
\[
 D=J_n-J_o=Y_C\Gamma_C^{-1}N_C-Y_A\Gamma_A^{-1}N_A.
\]
两个项都包含材料激发、内部反馈和接收返回，因而替换后的影响不是两个裸接收强度的差。如果共同核心不可逆，但两个端点本身稳定，就直接计算端点；不能把快速公式不能用，误写成物理空间不可行。

\subsection{材料信息描述响应，残差拉回描述当前任务}
在共同的实材料坐标中，定义数据曲率（data curvature）
\[
 \mathsf I_U=J_U^TJ_U.
\]
给一个材料变化 $v$，$v^T\mathsf I_Uv=\|J_Uv\|^2$，描述模型认为这种变化会产生多强的接收响应。这里用直立的 $\mathsf I_U$ 表示信息曲率，与前文双向尾部中的注入残差 $\mathcal I_U$ 区分。

材料步还需要常数梯度项（constant gradient term）
\[
 h_U=J_U^Tr+\ell,\qquad H_Us_U=-h_U.
\]
它保留接收响应与当前残差的方向、符号和复相位的实拉回。保存的数据把这个正号量写成 $b$；真正的法方程右端是 $-h_U$，不是 $h_U$。

令 $D=J_n-J_o$。直接展开两端点，可得
\[
 \Delta\mathsf I=J_o^TD+D^TJ_o+D^TD,
 \qquad \Delta h=D^Tr.
\]
再从 $H_ns_n=-h_n$ 减去旧端点的关系：
\[
 H_n(s_n-s_o)=-(h_n-h_o)-(H_n-H_o)s_o.
\]
由于正则项共同不变，$H_n-H_o=\Delta\mathsf I$，因此
\[
 \boxed{d=s_n-s_o=-H_n^{-1}(\Delta h+\Delta\mathsf I s_o).}
\]
这几行是从物理导数变化到实际更新的完整桥梁。只给出信息谱，会丢掉 $\Delta h$；只讨论残差右端，又会丢掉曲率改变。

一个简单代数例子能说明这一点。取 $r=1$、$\ell=0$、$\Lambda=1$。$J=1$ 与 $J=-1$ 的信息曲率都是1，却分别给出 $s=-1/2$ 与 $s=1/2$。它们有相同的信息量汇总，更新方向完全相反。这个例子是解释公式，不是新增电磁数据。

\subsection{信息增加和模型更准确要分别衡量}
近似模型可能过度放大某些材料响应，因此更大的信息体积不一定更接近完整物理。本文使用固定尺度 $P$，报告有效维数（effective dimension）与信息体积诊断（information-volume diagnostic）：
\[
 d_{\rm eff}=\sum_i\frac{\nu_i}{1+\nu_i},\qquad
 \mathcal D=\frac12\sum_i\log(1+\nu_i),
\]
其中 $\nu_i$ 是 $P^{-1/2}\mathsf I_UP^{-1/2}$ 的特征值。同时报告相对完整参考的信息保真度（information fidelity），而不是只奖励 $\mathcal D$ 变大。

实际比较把 $P=10^{-5}I$ 的先验尺度与Levenberg--Marquardt阻尼（LM damping）分开；局部求解仍使用原有总曲率。接收数据沿用固定数值归一化，尚无独立标定的噪声协方差，所以这些量作为尺度归一化曲率诊断，不宣称为硬件测量的互信息。较小的Gaussian材料空间可以算完整谱；voxel空间只在共同的16个材料方向上诊断，不能据此说完整空间没有弱方向。

\subsection{实际位移在完整二次目标上是否值得}
生成 $d$ 以后，用同一个完整冻结二次目标评价有符号收益（finite signed gain）：
\[
 Q=\Phi_F(x)-\Phi_F(x+d)=-g_F(x)^Td-\tfrac12d^TH_Fd.
\]
第一项检查位移是否朝向下降，第二项扣除这段有限位移的曲率代价。$Q>0$表示更忠实于当前完整GN更新；它没有把材料真值加入目标，所以不能直接保证图像更准确。

为检查这个标量，不必再计算一次完整伴随（full adjoint）。先缓存 $u=r+J_Fx$，对候选只算定向输出（directional output）$z=J_Fd$：
\[
 Q=-u^Tz-(\ell+\Lambda x)^Td
     -\tfrac12(\|z\|^2+d^T\Lambda d).
\]
复表示下第一项取 $\operatorname{Re}\langle u,z\rangle$。每个材料方向仍要支付所有照明的全波右端；这只是减少不必要的伴随计算，并不意味着免费。接受后更新空间、材料步及缓存输出，下一轮重算条件评分（conditional scoring）。

\subsection{高精度复核与严格证书的区别}
如果输出有经过验证的误差界 $e_u,e_z$，则收益误差可由
\[
 \epsilon_Q=e_u\|\widehat z\|+(\|\widehat u\|+\|\widehat z\|)e_z+e_ue_z+\tfrac12e_z^2
\]
控制。接受一个动作需要其收益下界为正；证明整个声明邻域没有好动作，需要每个可行动作的收益上界都低于容差。这是两个不同判断。

现有全波输出尚无这样闭合的严格界，所以实验采用双精度定向复核（high-fidelity directional check），不把数值容差当作确定性证书（deterministic certificate）。若只检查少量候选，没有找到好动作，就保留未解决状态；三次动作预算用完也不是局部最优的证明。

\subsection{怎样阅读方向分解图}
信息方向固定为同一个完整参考的特征子空间，而不是每换一个模型就重新给特征向量编号。Gaussian空间中，总求解曲率在这些方向上的权重是数据曲率加先验和LM项，因此各带的步误差能量可以精确相加；参考步尚有微小梯度时，另保留其线性交叉项。

一个约化模型认为很弱的方向，可能只是它没有表示出来；完整模型仍可能对此很敏感。这叫约化模型假弱（ROM-false-weak），与真实数据弱（true data-weak）不同。成功交换可能修复假弱、减少过度曲率或调整残差方向，并不必然是在给真实弱方向补先验。
